\chapter{Stellar Photometry: Flux, Magnitudes, and Radii}

Chapter~3 used photometric distance indicators (spectroscopic parallax, standard candles) without yet justifying the underlying flux measurements themselves. Photometry is the quantitative measurement of electromagnetic flux received from celestial objects. From the measured radiant flux, astronomers deduce fundamental astrophysical properties of stars, including luminosity ($L$), effective surface temperature ($T_{\rm eff}$), radius ($R_*$), and distance ($d$). In this chapter, we develop the inverse-square law, Pogson's logarithmic magnitude system, blackbody radiation thermodynamics, and the determination of stellar physical dimensions.

\section{Radiant Energy, Flux, and the Inverse-Square Law}

\begin{definition}[Luminosity and Flux]
	The \emph{luminosity} $L$ of a star is the total energy radiated into space per unit time (power in Watts, $\mathrm{W}$). The \emph{radiant flux} $F$ (or apparent brightness) is the energy received per unit time per unit collecting area perpendicular to the line of sight:
	\begin{keyresult}
		\begin{equation}
			F = \frac{L}{4\pi d^2}\quad \left[\mathrm{W\cdot m^{-2}}\right],
			\label{eq:inverse_square_flux}
		\end{equation}
	\end{keyresult}
	where $d$ is the distance from the source to the observer.
\end{definition}

Eq.~\eqref{eq:inverse_square_flux} is the geometric \emph{inverse-square law of radiation}, arising from energy conservation as radiant power spreads uniformly over a concentric spherical surface of area $4\pi d^2$ in an isotropic, non-absorbing medium.

\begin{example}[Solar Constant]
	The Sun has luminosity $L_\odot = 3.828 \times 10^{26}\,\mathrm{W}$. At the mean Earth-Sun distance ($1\,\mathrm{AU} = 1.4959787 \times 10^{11}\,\mathrm{m}$), the total solar irradiance (the \emph{solar constant}) at the top of Earth's atmosphere is:
	\[
	F_\odot = \frac{3.828 \times 10^{26}}{4\pi (1.4959787 \times 10^{11})^2} \approx 1361\,\mathrm{W\cdot m^{-2}}.
	\]
\end{example}

\section{The Astronomical Magnitude System}

Human vision responds logarithmically to perceived optical brightness (the Weber-Fechner law). In the 2nd century BCE, Hipparchus cataloged stars into six discrete brightness classes ($1^{\rm st}$ magnitude for the brightest down to $6^{\rm th}$ magnitude for the faintest naked-eye stars).

\subsection{Apparent Magnitude and Pogson's Formula}

In 1856, Norman Pogson formalized this ancient scale mathematically after noticing that a difference of 5 magnitudes corresponds to a flux ratio of exactly $100:1$. Thus, a single magnitude step corresponds to a flux ratio of:
\begin{equation}
	\sqrt[5]{100} = 100^{1/5} = 10^{0.4} \approx 2.511886.
\end{equation}

\begin{definition}[Pogson's Formula for Apparent Magnitude]
	The difference in \emph{apparent magnitude} $m_1 - m_2$ between two sources with measured fluxes $F_1$ and $F_2$ is:
	\begin{keyresult}
		\begin{equation}
			m_1 - m_2 = -2.5\log_{10}\left(\frac{F_1}{F_2}\right).
			\label{eq:pogson_difference}
		\end{equation}
	\end{keyresult}
	Inverting Pogson's relation:
	\begin{equation}
		\frac{F_1}{F_2} = 10^{-0.4(m_1 - m_2)}.
		\label{eq:flux_ratio_from_mag}
	\end{equation}
\end{definition}

On an absolute scale, apparent magnitude is defined relative to a standard zero-point flux $F_0$:
\begin{equation}
	m = -2.5\log_{10}\left(\frac{F}{F_0}\right) = -2.5\log_{10}F + C_0.
	\label{eq:apparent_mag_single}
\end{equation}

\subsection{Absolute Magnitude and Distance Modulus}

Because apparent magnitude $m$ depends on both intrinsic luminosity $L$ and distance $d$, it cannot directly measure a star's true physical power.

\begin{definition}[Absolute Magnitude]
	The \emph{absolute magnitude} $M$ is defined as the apparent magnitude a star would possess if it were located at a standard fiducial distance of exactly $d_0 = 10\,\mathrm{pc}$.
\end{definition}

Applying Pogson's formula between a star's actual distance $d$ and the fiducial distance $10\,\mathrm{pc}$:
\[
m - M = -2.5\log_{10}\left(\frac{F(d)}{F(10\,\mathrm{pc})}\right) = -2.5\log_{10}\left(\frac{L/(4\pi d^2)}{L/(4\pi (10)^2)}\right) = -2.5\log_{10}\left(\frac{10}{d}\right)^2 = 5\log_{10}\left(\frac{d}{10}\right).
\]
\begin{keyresult}
	\begin{equation}
		\mu \equiv m - M = 5\log_{10}\left(\frac{d}{10\,\mathrm{pc}}\right) = 5\log_{10}d\,[\mathrm{pc}] - 5.
		\label{eq:distance_modulus}
	\end{equation}
\end{keyresult}
The quantity $\mu = m - M$ is known as the \emph{distance modulus}. In terms of trigonometric parallax $p$ (in arcseconds):
\begin{equation}
	m - M = -5\log_{10}p - 5.
	\label{eq:distance_modulus_parallax}
\end{equation}

\subsection{Photometric Passbands, Color Indices, and Interstellar Extinction}

Astronomical detectors measure flux through calibrated optical transmission filters (e.g., the Johnson-Cousins $UBVRI$ system or SDSS $ugriz$ system):
\begin{itemize}
	\item $U$ (Ultraviolet, $\lambda_{\rm eff} \approx 365\,\mathrm{nm}$)
	\item $B$ (Blue, $\lambda_{\rm eff} \approx 440\,\mathrm{nm}$)
	\item $V$ (Visual, $\lambda_{\rm eff} \approx 550\,\mathrm{nm}$)
	\item $R$ (Red, $\lambda_{\rm eff} \approx 640\,\mathrm{nm}$)
	\item $I$ (Infrared, $\lambda_{\rm eff} \approx 790\,\mathrm{nm}$)
\end{itemize}

\begin{definition}[Color Index]
	A \emph{color index} is the difference in apparent magnitude between two distinct wavelength passbands (e.g., $B - V$ or $U - B$). Because flux appears in the denominator of the logarithmic ratio:
	\begin{equation}
		B - V = -2.5\log_{10}\left(\frac{F_B}{F_V}\right) + \text{const}.
		\label{eq:color_index_def}
	\end{equation}
	Hotter stars emit relatively more blue flux ($F_B > F_V$), resulting in \emph{smaller (or negative) color indices} ($B - V < 0$). Cooler stars have $F_B < F_V$, yielding \emph{larger positive color indices} ($B - V > 0$).
\end{definition}

\paragraph{Interstellar Reddening and Extinction.}
Interstellar dust grains preferentially scatter shorter-wavelength blue photons (Rayleigh-like scattering $\sigma_{\rm dust} \propto \lambda^{-\gamma}$ with $\gamma \sim 1$; the microphysical origin of this power law in Mie and Rayleigh grain scattering is derived in Chapter~8). This dims and reddens the transmitted starlight. The observed distance modulus is modified by the total visual extinction $A_V$:
\begin{equation}
	m_V - M_V = 5\log_{10}d - 5 + A_V.
	\label{eq:extinction_modulus}
\end{equation}
The \emph{color excess} $E(B-V)$ is defined as:
\begin{equation}
	E(B-V) \equiv (B - V)_{\rm observed} - (B - V)_{\rm intrinsic}.
	\label{eq:color_excess}
\end{equation}
In the diffuse interstellar medium of the Milky Way, the ratio of total to selective extinction is standardly:
\begin{keyresult}
	\begin{equation}
		R_V \equiv \frac{A_V}{E(B-V)} \approx 3.1.
		\label{eq:rv_extinction}
	\end{equation}
\end{keyresult}

\section{Thermal Radiation and the Stefan-Boltzmann Law}

To relate radiant emission to stellar physical properties, we model stellar photospheres as thermodynamic blackbodies.

\subsection{Planck's Law of Blackbody Radiation}

A blackbody in thermodynamic equilibrium at temperature $T$ radiates a specific intensity $B_\lambda(T)$ (energy per unit time, area, solid angle, and wavelength):
\begin{keyresult}
	\begin{equation}
		B_\lambda(T) = \frac{2hc^2}{\lambda^5} \frac{1}{\exp\!\left(\dfrac{hc}{\lambda k_B T}\right) - 1}\quad \left[\mathrm{W\cdot m^{-2}\cdot sr^{-1}\cdot m^{-1}}\right],
		\label{eq:planck_wavelength}
	\end{equation}
\end{keyresult}
where $h = 6.62607 \times 10^{-34}\,\mathrm{J\cdot s}$ is Planck's constant, $c = 2.99792 \times 10^8\,\mathrm{m/s}$ is the speed of light, and $k_B = 1.38065 \times 10^{-23}\,\mathrm{J/K}$ is Boltzmann's constant.

In frequency space:
\begin{equation}
	B_\nu(T) = \frac{2h\nu^3}{c^2} \frac{1}{\exp\!\left(\dfrac{h\nu}{k_B T}\right) - 1}\quad \left[\mathrm{W\cdot m^{-2}\cdot sr^{-1}\cdot Hz^{-1}}\right].
	\label{eq:planck_frequency}
\end{equation}

\subsection{Wien's Displacement Law}

To find the wavelength of peak emission, write $B_\lambda \propto \lambda^{-5}(e^x-1)^{-1}$ with $x \equiv hc/(\lambda k_BT)$, and note that since $x \propto 1/\lambda$, the chain rule gives $dx/d\lambda = -x/\lambda$. Differentiating with respect to $\lambda$ (product and chain rule) and setting the result to zero:
\begin{equation}
	\dv{B_\lambda}{\lambda} \propto -5\lambda^{-6}(e^x-1)^{-1} + \lambda^{-5}(e^x-1)^{-2}e^x\left(\frac{x}{\lambda}\right) = 0.
\end{equation}
Multiplying through by $-\lambda^{6}(e^x-1)$ gives the transcendental condition on $x$:
\begin{equation}
	\frac{x e^x}{e^x - 1} = 5, \quad \text{where } x \equiv \frac{hc}{\lambda_{\max} k_B T} \approx 4.965114 \quad \text{(solved numerically).}
\end{equation}
This establishes \emph{Wien's Displacement Law}:
\begin{keyresult}
	\begin{equation}
		\lambda_{\max} T = \frac{hc}{4.965114\,k_B} \approx 2.89777 \times 10^{-3}\,\mathrm{m\cdot K} = 2898\,\mathrm{\mu m\cdot K}.
		\label{eq:wien_law}
	\end{equation}
\end{keyresult}

\subsection{The Stefan-Boltzmann Law and Effective Temperature}

Integrating Planck's specific intensity over all wavelengths and across the outward hemisphere ($\int_0^{2\pi}d\phi\int_0^{\pi/2}\cos\theta\sin\theta\,d\theta = \pi$, the projected solid angle of a hemisphere) gives the total emergent surface flux $F_{\rm surface} = \pi\int_0^\infty B_\lambda(T)\,d\lambda$. Substituting $x = hc/(\lambda k_BT)$ (so $\lambda = hc/(xk_BT)$ and $d\lambda = -hc\,dx/(xk_BT)^2$... more directly, in frequency form with $x=h\nu/k_BT$) converts the integral into a pure number times $T^4$:
\begin{equation}
	\int_0^\infty B_\nu(T)\,d\nu = \frac{2k_B^4T^4}{h^3c^2}\int_0^\infty \frac{x^3}{e^x-1}\,dx = \frac{2k_B^4T^4}{h^3c^2}\cdot\frac{\pi^4}{15},
\end{equation}
using the standard result $\int_0^\infty x^3/(e^x-1)\,dx = \pi^4/15$ (a tabulated integral, obtained by expanding $1/(e^x-1)=\sum_{n=1}^\infty e^{-nx}$ and summing the resulting Gamma-function integrals, $\sum_n \Gamma(4)/n^4 = 6\zeta(4) = 6\cdot\pi^4/90 = \pi^4/15$). Multiplying by $\pi$ gives
\begin{equation}
	F_{\rm surface} = \sigma_{\rm SB} T^4,
	\label{eq:surface_flux_sb}
\end{equation}
where the \emph{Stefan-Boltzmann constant} collects all the remaining factors:
\begin{equation}
	\sigma_{\rm SB} = \frac{2\pi^5 k_B^4}{15 c^2 h^3} \approx 5.670374 \times 10^{-8}\,\mathrm{W\cdot m^{-2}\cdot K^{-4}}.
	\label{eq:stefan_boltzmann_constant}
\end{equation}

\begin{definition}[Effective Temperature]
	The \emph{effective surface temperature} $T_{\rm eff}$ of a star with total luminosity $L$ and radius $R_*$ is the temperature of an ideal blackbody having the identical surface area and total bolometric power:
	\begin{keyresult}
		\begin{equation}
			L = 4\pi R_*^2\,\sigma_{\rm SB} T_{\rm eff}^4.
			\label{eq:stefan_boltzmann_luminosity}
		\end{equation}
	\end{keyresult}
\end{definition}

\section{Determination of Stellar Physical Radii}

\subsection{Radius from Bolometric Luminosity and Temperature}

Normalizing Eq.~\eqref{eq:stefan_boltzmann_luminosity} relative to solar values ($L_\odot = 3.828 \times 10^{26}\,\mathrm{W}$, $R_\odot = 6.957 \times 10^8\,\mathrm{m}$, $T_{\rm eff,\odot} = 5772\,\mathrm{K}$):
\begin{keyresult}
	\begin{equation}
		\frac{R_*}{R_\odot} = \sqrt{\frac{L_*}{L_\odot}}\left(\frac{T_{\rm eff,\odot}}{T_{\rm eff}}\right)^2.
		\label{eq:stellar_radius_solar_units}
	\end{equation}
\end{keyresult}
In terms of absolute bolometric magnitude $M_{\rm bol}$:
\begin{equation}
	\log_{10}\left(\frac{R_*}{R_\odot}\right) = -0.2(M_{\rm bol,*} - M_{\rm bol,\odot}) - 2\log_{10}\left(\frac{T_{\rm eff}}{T_{\rm eff,\odot}}\right),
	\label{eq:radius_from_mbol}
\end{equation}
where $M_{\rm bol,\odot} = +4.74\,\mathrm{mag}$.

\subsection{Angular Diameter ($\theta$)}

A spherical star of physical radius $R_*$ located at distance $d$ subtends an angular diameter $\theta$:
\begin{equation}
	\tan\left(\frac{\theta}{2}\right) = \frac{R_*}{d} \quad\Longrightarrow\quad \theta \approx \frac{2 R_*}{d}\quad (\text{radians}).
	\label{eq:angular_diameter_rad}
\end{equation}
Converting $\theta$ to arcseconds ($1\,\mathrm{rad} = 206{,}264.8''$):
\begin{keyresult}
	\begin{equation}
		\theta\,[\mathrm{mas}] = 206{,}264{,}800 \times \frac{2 R_*}{d} = 9.300 \times 10^{-3} \left(\frac{R_*}{R_\odot}\right)\left(\frac{1\,\mathrm{pc}}{d}\right).
		\label{eq:angular_diameter_mas}
	\end{equation}
\end{keyresult}

\begin{example}[Angular Diameter of Betelgeuse]
	Betelgeuse ($\alpha$ Orionis) is a red supergiant with $R_* \approx 764\,R_\odot$ at a distance $d \approx 168\,\mathrm{pc}$ ($p \approx 5.95\,\mathrm{mas}$). Its expected angular diameter is:
	\[
	\theta = 9.300 \times 10^{-3} \times \frac{764}{168} \approx 42.3\,\mathrm{mas}.
	\]
	This matches direct optical interferometric measurements ($42\text{--}44\,\mathrm{mas}$), demonstrating that giant stars can be directly resolved with modern optical/IR interferometers (such as the VLTI).
\end{example}

Resolving an angular diameter of tens of milliarcseconds, as in the Betelgeuse example above, is only possible because of the interferometric and diffraction-limited resolution machinery developed next in Chapter~5.
